\chapter{Le articolazioni dell'arto superiore}

% ---------------------------------------------------------------------------
\section{Generalità e cinto scapolare}

\textbf{Scheletro appendicolare}: ossa degli arti + cinti (cingoli) che li collegano al tronco
$\rightarrow$ \textbf{cinto scapolare} (toracico, arti superiori) e \textbf{cinto pelvico} (arti
inferiori).
\begin{itemize}
  \item arto superiore: cinto scapolare (clavicola + scapola) + arto libero = braccio (omero),
    avambraccio (ulna + radio), mano (carpo, metacarpo, falangi);
  \item articolazioni: sterno-clavicolare, acromio-clavicolare, glenomerale, gomito, radio-ulnare
    distale, articolazioni della mano;
  \item il cinto scapolare (2 clavicole avanti, 2 scapole dietro) \textbf{non} si articola con la
    colonna; comprende sternoclavicolare, acromioclavicolare, scapolotoracica (scorrimento
    funzionale) e, con l'omero, la scapolo-omerale (gleno-omerale).
\end{itemize}

% ---------------------------------------------------------------------------
\section{Articolazione sterno-clavicolare}

Unisce l'estremità sternale della clavicola all'incisura clavicolare del manubrio. Tre legamenti:
\begin{itemize}
  \item \textbf{inter-clavicolare}: collega le estremità sternali delle due clavicole (incisura
    giugulare del manubrio);
  \item \textbf{costo-clavicolare}: cono rovesciato, fascio anteriore e posteriore;
  \item \textbf{sterno-clavicolare}: banda anteriore (fino alla I cartilagine costale) e posteriore
    (meno robusta).
\end{itemize}
Disco articolare interposto; in rapporto con la I costa. Movimenti: protrazione/retrazione,
elevazione/depressione della clavicola (rotolamento + scivolamento, \textit{roll} e \textit{slide}).

% ---------------------------------------------------------------------------
\section{Articolazione acromion-clavicolare}

Unisce l'estremità acromiale della clavicola all'acromion; soggetta a traumi e lussazioni.
\begin{itemize}
  \item mezzi di unione: \textbf{acromioclavicolare} + \textbf{coracoclavicolare} (fasci trapezoide
    e conoide);
  \item \textbf{coracoacromiale}: con acromion, processo coracoideo e coracoclavicolare forma
    l'\textbf{arco coracoacromiale} (sopra la testa dell'omero);
  \item movimenti: piccoli scivolamenti (variano l'angolo scapola-clavicola) e rotazione della
    scapola (interna/esterna, \textit{tilting} anteriore/posteriore, verso alto/basso).
\end{itemize}

\subsection{Lesioni dell'acromioclavicolare}
Caduta sulla spalla/braccio esteso $\rightarrow$ lussazione + lesione dei legamenti. Tre gradi:
\begin{enumerate}
  \item stiramento dei legamenti;
  \item rottura del legamento acromioclavicolare;
  \item lussazione totale (rottura di coracoclavicolare + acromioclavicolare).
\end{enumerate}
Nella lussazione: rilievo della clavicola con sollevamento della pelle; premendo si riduce ma la
deformità si ripresenta (\textbf{segno del tasto di pianoforte}). Cause: trauma diretto, caduta sul
moncone, traumi sportivi. Segni: rilievo della clavicola, dolore localizzato. Trattamento:
mobilizzazioni + terapia manuale $\rightarrow$ esercizio terapeutico (nelle prime fasi sotto i 90°
di elevazione) e rinforzo $\rightarrow$ eventuale terapia strumentale.

\subsection{Legamenti propri della scapola}
Origine e inserzione solo sulla scapola: coracoacromiale, coracoclavicolare, trasverso superiore.
Alla cavità glenoidea: labbro glenoideo (ne approfondisce il profilo), tendine del bicipite, capsula,
legamenti glenomerali superiore/medio/inferiore.

% ---------------------------------------------------------------------------
\section{Articolazione glenomerale}

Tra cavità glenoidea e testa dell'omero. Capi emisferici $\rightarrow$ \textbf{enartrosi} (a sfera),
3 assi, 6 movimenti; il maggior \textbf{range} di tutto il corpo.
\begin{itemize}
  \item capsula + legamenti glenomerali (in particolare il \textbf{coracoomerale}) tengono insieme
    le superfici;
  \item \textbf{membrana sinoviale}: riveste la capsula e la parte articolare dell'osso (a contatto
    con la membrana fibrosa fuori, col liquido sinoviale dentro); associati: borsa sottocoracoidea,
    legamento trasverso dell'omero, tendine del capo lungo del bicipite (nel solco bicipitale, con
    guaina sinoviale), borsa sottotendinea del sottoscapolare.
\end{itemize}

\subsection{Cinematica gleno-omerale}
Quattro gradi di libertà:
\begin{itemize}
  \item sagittale: flessione 0--180°, estensione 0--45°;
  \item frontale: abduzione 0--120°, adduzione 0--30/45°;
  \item trasversale: flessione orizzontale 0--140°, estensione orizzontale 0--30°;
  \item orizzontale: rotazione interna 0--70°, esterna 0--80°.
\end{itemize}

% ---------------------------------------------------------------------------
\section{Le articolazioni del gomito}

Tre articolazioni (omero-ulnare, omero-radiale, radio-ulnare prossimale) tra estremità distale
dell'omero e prossimali di radio/ulna; funzionano da cardine (ginglimo). Stabilità principale =
incastro troclea omerale/incisura trocleare dell'ulna.

\subsection{Omero-ulnare}
Troclea dell'omero nell'incisura trocleare dell'ulna $\rightarrow$ \textbf{ginglimo angolare} (a
troclea, 1 asse, 2 movimenti); flesso-estensione dell'avambraccio.

\subsection{Omero-radiale}
Condilo dell'omero con la testa radiale (capitello) $\rightarrow$ \textbf{condilartrosi}.

\subsection{Radio-ulnare prossimale}
Capitello del radio con l'incisura radiale dell'ulna $\rightarrow$ \textbf{ginglimo laterale}
(trocoide, 1 asse); prono-supinazione.

\subsection{Stabilità e traumatologia}
Molto stabile (previene movimenti laterali/rotatori); capsula spessa con legamenti robusti:
\begin{itemize}
  \item \textbf{collaterale ulnare} (mediale): epicondilo mediale $\rightarrow$ margine mediale
    dell'incisura trocleare;
  \item \textbf{collaterale radiale} (laterale): epicondilo laterale $\rightarrow$ 3 fasci (anteriore,
    medio, posteriore);
  \item \textbf{anulare del radio}: circonda la testa radiale, fissato all'incisura radiale;
    \textbf{legamento quadrato}: collo del radio $\rightarrow$ incisura radiale;
  \item \textbf{membrana interossea}: tra radio e ulna, separa i muscoli anteriori dai posteriori.
\end{itemize}
\textbf{Fratture del gomito}: 7\% di tutte le fratture $\rightarrow$ 33\% distale dell'omero, 10\%
olecrano, 35\% capitello radiale; traumi indiretti (caduta in estensione/flessione).

\subsection{Nota clinica: sublussazione del capitello radiale}
Nei bambini, per trazione sull'avambraccio (braccio tirato in alto con avambraccio intraruotato)
$\rightarrow$ strappo del legamento anulare (lasso sul collo del radio) $\rightarrow$ rifiuto a
muovere il gomito. Riduzione: extrarotazione dell'avambraccio + flessione del gomito.

% ---------------------------------------------------------------------------
\section{Articolazione radio-ulnare distale}

Tra capitello dell'ulna e incisura ulnare del radio; \textbf{trocoide} (ginglimo laterale). Con la
radio-ulnare prossimale $\rightarrow$ prono-supinazione.

% ---------------------------------------------------------------------------
\section{Le articolazioni della mano}

Radio-carpica, inter-carpiche, carpo-metacarpiche, inter-metacarpiche, metacarpo-falangee,
inter-falangee.
\begin{itemize}
  \item \textbf{radio-carpica}: superficie distale del radio (+ disco) con la fila prossimale del
    carpo $\rightarrow$ \textbf{condiloartrosi}; stabilizzata da capsula e collaterale ulnare del
    carpo.
\end{itemize}

\subsection{Inter-carpiche}
\begin{itemize}
  \item fila prossimale: \textbf{artrodie} (scafoide, semilunare, piramidale); il pisiforme si
    articola con la faccia palmare del piramidale;
  \item fila distale: \textbf{artrodie} (trapezio, trapezoide, capitato, uncinato);
  \item \textbf{mediocarpica} (tra le due file, escluso il pisiforme): artrodia + condiloidea.
\end{itemize}
Movimenti della mano: flessione, estensione, abduzione (deviazione radiale), adduzione (deviazione
ulnare), circumduzione.

\subsection{Carpo-metacarpiche e delle dita}
\begin{itemize}
  \item \textbf{carpo-metacarpiche}: artrodie (eccetto il pollice, \textbf{a sella}); flessione,
    estensione, modesta lateralità; il I metacarpale è il più mobile;
  \item \textbf{inter-metacarpiche}: artrodie (basi degli ultimi 4 metacarpali), solo scivolamento;
  \item \textbf{metacarpo-falangee}: a condilo; flesso-estensione + modeste abduzione/adduzione/
    circumduzione;
  \item \textbf{inter-falangee}: a troclea; due per dito (una sola per il pollice); ampia
    flesso-estensione.
\end{itemize}
